%% kapa: exam=1 round=34 year=2023 period=1 form=A subject=economics subject_ko=경제학원론 seq=36 number=76
%% source: Q-Net 감정평가사 기출문제 (cst00309) artlSeq=5211789 file=raw/1008/5211789/2023년도 제34회 감정평가사 1차시험 1교시.pdf
%% conversion: kordoc 4.24.51 -> tools/gen1.py (auto); review=proofread-sonnet; flags=none
\begin{question}{76}
총수요곡선 및 총공급곡선에 관한 설명으로 옳지 않은 것을 모두 고른 것은?
\begin{framedtext}
ㄱ. 총수요곡선은 물가수준과 재화 및 용역의 수요량간의 관계를 보여준다. \par
ㄴ. 통화수요 또는 투자가 이자율에 영향을 받지 않을 경우 총수요곡선은 수평이 된다. \par
ㄷ. 단기적으로 가격이 고정되어 있을 경우 총공급곡선은 수평이 된다. \par
ㄹ. 정부지출의 변화는 총수요곡선 상에서의 변화를 가져온다. \par
\end{framedtext}
\begin{choices}
  \item ㄱ, ㄴ
  \item ㄱ, ㄷ
  \item ㄴ, ㄷ
  \item ㄴ, ㄹ
  \item ㄷ, ㄹ
\end{choices}
\end{question}
\answer{4}
